\documentclass[12pt]{article}
\usepackage[margin=1in]{geometry}
\usepackage[all]{xy}


\usepackage{amsmath,amsthm,amssymb,color,latexsym}
\usepackage{geometry}        
\geometry{letterpaper}  
\usepackage{boxedminipage}
\usepackage{graphicx}
\usepackage{hyperref}
\hypersetup{
    colorlinks=true,
    citecolor=blue,
    linkcolor=blue,
    filecolor=cyan,
    urlcolor=magenta,
}
\usepackage[capitalize,nameinlink]{cleveref}

\newtheorem{problem}{Problem}
\newtheorem{lemma}{Lemma}


\newenvironment{solution}[1][\it{Solution}]{\textbf{#1. } }{$\square$}
\newcommand{\pointsbox}[1]{%
  \colorbox{red!70}{\strut\textcolor{white}{\sffamily\bfseries\ \ #1\ \ }}%
}
\newcommand{\pagetag}[1]{%
  \colorbox{red!70}{\strut\textcolor{white}{\sffamily\bfseries\ \ #1\ \ }}%
}
\newcommand{\striptitle}[2]{%
  \noindent
  \begin{minipage}{0.07\linewidth}\centering \pagetag{#1}\end{minipage}%
  \hfill
  \begin{minipage}{0.89\linewidth}
    {\large\scshape #2}
  \end{minipage}\par\vspace{0.75em}
}
% --- Environments ---
\theoremstyle{plain}
\newtheorem{theorem}{Theorem}
\newtheorem{corollary}[theorem]{Corollary}
\theoremstyle{definition}
\newtheorem{definition}[theorem]{Definition}
\newtheorem{remark}[theorem]{Remark}

\newenvironment{boxedalgo}
  {\begin{center}\begin{boxedminipage}{1\linewidth}}
  {\end{boxedminipage}\end{center}}

\newenvironment{inlinealgo}[1]
  {\noindent\hrulefill\\\noindent{\bf {#1}}}
  {\hrulefill}


% --- Shortcuts ---
\newcommand{\R}{\mathbb{R}}
\newcommand{\inner}[2]{\left\langle #1,#2 \right\rangle}
\newcommand{\proj}{\mathrm{proj}}
\newcommand{\Id}{\mathbf{I}}
\newcommand{\ot}{\otimes}
\newcommand{\mc}{\mathcal}

\begin{document}
\noindent CS 578 Statistical Machine Learning (Fall 2026)\hfill Problem Set \# 5\\
\{\{Your name\}\} \hfill \textbf{Due: November 2nd}

\noindent\hrulefill

\subsection*{Problem 1. Markov Properties.}
Let \(G=([n],E)\) be a simple undirected graph, and let \(X=(X_1,\ldots,X_n)\) have a joint distribution \(\mu\), not necessarily strictly positive, on finite state spaces. For \(A\subseteq[n]\), write \(X_A=(X_i)_{i\in A}\), and let \(\partial_i=\{j\in[n]:\{i,j\}\in E\}\) denote the neighbors of vertex \(i\). Let \(U\perp V\mid W\) denote that \(U\) and \(V\) are conditionally independent given \(W\).

Recall the following definitions from lecture:
\begin{itemize}
    \item \textbf{Global Markov property.} For any pairwise disjoint sets \(A,B,S\subseteq[n]\), if every path from a vertex in \(A\) to a vertex in \(B\) passes through \(S\), then $X_A\perp X_B\mid X_S$.

    \item \textbf{Local Markov property.} For every \(i\in[n]\), conditioned on its neighbors, \(X_i\) is independent of all remaining variables: $X_i\perp X_{[n]\setminus(\{i\}\cup\partial_i)}
        \mid X_{\partial_i}$.
        
    \item \textbf{Pairwise Markov property.} For every pair of distinct, nonadjacent vertices \(i,j\in[n]\), $X_i\perp X_j\mid X_{[n]\setminus\{i,j\}}$.
\end{itemize}

\begin{enumerate}
    \item Prove that the global Markov property implies the local Markov property.
    \item Prove that the local Markov property implies the pairwise Markov property.

    \emph{Hint: you may first show that $U\perp(V,W)\mid Z\Longrightarrow U\perp V\mid(Z,W)$.}
\end{enumerate}

%You may not assume that \(\mu\) is strictly positive. Conditional probabilities need only be considered when the conditioning event has positive probability.

\vspace{1em}
\noindent{\bfseries Solution.}\newline 
\begin{enumerate}
    \item Fix a vertex \(i\), and take
    \[
    A=\{i\},\qquad S=\partial_i,\qquad
    B=[n]\setminus(\{i\}\cup\partial_i).
    \]
    These sets are pairwise disjoint. Every path from \(i\) to a
    vertex of \(B\) must first visit a neighbor of \(i\), hence must
    pass through \(S\). The global Markov property therefore gives
    \(X_i\perp X_B\mid X_{\partial_i}\), which is precisely the
    local Markov property.

    \item First prove the weak union implication in the hint. Suppose
    \(U\perp(V,W)\mid Z\), and fix \((z,w)\) with
    \(\Pr(Z=z,W=w)>0\). Conditional independence gives
    \begin{align*}
    \Pr(U=u,V=v\mid Z=z,W=w)
    &=\frac{\Pr(U=u,V=v,W=w\mid Z=z)}
            {\Pr(W=w\mid Z=z)}\\
    &=\Pr(U=u\mid Z=z)\Pr(V=v\mid Z=z,W=w).
    \end{align*}
    Summing over \(v\) shows that
    \(\Pr(U=u\mid Z=z,W=w)=\Pr(U=u\mid Z=z)\).
    Substituting this into the factorization above proves
    \(U\perp V\mid(Z,W)\).

    Now let \(i,j\) be distinct nonadjacent vertices and set
    \(R=[n]\setminus(\{i,j\}\cup\partial_i)\).
    Since \(j\notin\partial_i\), the local Markov property at \(i\)
    is
    \[
    X_i\perp(X_j,X_R)\mid X_{\partial_i}.
    \]
    Since \(\partial_i\cup R=[n]\setminus\{i,j\}\), weak union yields
    \[
    X_i\perp X_j\mid X_{[n]\setminus\{i,j\}}.
    \]
    This is the pairwise Markov property. All conditional probabilities
    used above condition on positive-probability events, so strict
    positivity of \(\mu\) is unnecessary.
\end{enumerate}
\newpage

\subsection*{Problem 2. Higher-order Tensor Decomposition.}
Let \(u \odot v\) denote the Khatri--Rao product of two vectors:
if \(u \in \mathbb{R}^m\) and \(v \in \mathbb{R}^n\), then
\[
u \odot v = \operatorname{vec}(uv^\top) \in \mathbb{R}^{mn},
\]
where \(\operatorname{vec}\) flattens a matrix column by column.

The \emph{Kruskal rank} (or \(k\)-rank) of a collection
of vectors \(u_1,u_2,\ldots,u_m \in \mathbb{R}^n\) is the largest
\(k\in \{0,\dots,m\}\) such that every subset of \(k\) vectors is linearly independent.

In this problem, we explore properties of the Khatri--Rao product
and use them to design algorithms for decomposing higher-order tensors.

\begin{enumerate}
    \item[(a)]
    Let \(k_u\) and \(k_v\) be the \(k\)-ranks of
    \(u_1,u_2,\ldots,u_m\) and \(v_1,v_2,\ldots,v_m\), respectively. Assume all $u_i$ and $v_i$ are nonzero.
    Prove that the \(k\)-rank of
    \[
    u_1 \odot v_1,\ u_2 \odot v_2,\ \ldots,\ u_m \odot v_m
    \]
    is at least \(\min(k_u+k_v-1,m)\).

    \item[(b)]
    Construct a family of examples where the \(k\)-rank of
    \[
    u_1 \odot u_1,\ u_2 \odot u_2,\ \ldots,\ u_m \odot u_m
    \]
    is exactly \(2k_u-1\), and not any larger.
    To make this nontrivial, your examples must satisfy
    \(m>2k_u-1\).

    \item[(c)]
    Given an \(n \times n \times n \times n \times n\)
    fifth-order tensor
    \[
    T = \sum_{i=1}^{r} a_i^{\otimes 5},
    \]
    give an algorithm for finding its factors that works for
    \(r=2n-1\), under appropriate conditions on the factors
    \(a_1,a_2,\ldots,a_r \in \mathbb{R}^n\).

    \emph{Hint: Reduce to the third-order case.}
\end{enumerate}

\noindent{\bfseries Solution.}\newline 
\begin{enumerate}
    \item[(a)] Set \(K=\min(k_u+k_v-1,m)\). Since all factors are
    nonzero, \(k_u,k_v\geq1\). Fix any \(S\subseteq[m]\) with
    \(|S|=s\leq K\), and suppose
    \[
    \sum_{j\in S}c_j(u_j\odot v_j)=0,
    \qquad\text{equivalently}\qquad
    \sum_{j\in S}c_j u_jv_j^\top=0.
    \]
    Fix \(i\in S\). Choose \(J\subseteq S\setminus\{i\}\) with
    \(|J|=\min(k_u-1,s-1)\). The vectors indexed by
    \(J\cup\{i\}\) are linearly independent, so there is a linear
    functional \(\ell\) such that
    \(\ell(u_i)=1\) and \(\ell(u_j)=0\) for \(j\in J\).
    Applying \(\ell\) to the first factor of the matrix identity gives
    \[
    \sum_{j\in S\setminus J}c_j\ell(u_j)v_j=0.
    \]
    There are
    \[
    |S\setminus J|=\max(1,s-k_u+1)\leq k_v
    \]
    remaining vectors, so they are linearly independent. In particular,
    \(c_i\ell(u_i)=c_i=0\). Since \(i\) was arbitrary, all coefficients
    vanish. Thus every subset of at most \(K\) product vectors is
    independent, proving the claimed lower bound.

    \item[(b)] For any integer \(k\geq1\) and any \(m\geq2k\), choose
    distinct real numbers \(t_1,\ldots,t_m\), and define
    \[
    u_i=(1,t_i,t_i^2,\ldots,t_i^{k-1})^\top\in\mathbb R^k.
    \]
    Any \(k\) of these vectors form an invertible Vandermonde matrix,
    so \(k_u=k\). For \(0\leq d\leq2k-2\), let \(H_d\) be the
    \(k\times k\) matrix whose \((p,q)\)-entry is \(1\) when
    \(p+q-2=d\), and \(0\) otherwise. Then
    \[
    u_i\odot u_i
    =\sum_{d=0}^{2k-2}t_i^d\operatorname{vec}(H_d).
    \]
    Hence all product vectors lie in a subspace of dimension at most
    \(2k-1\). Because \(m\geq2k\), every set of \(2k\) of them is
    dependent, giving \(k\)-rank at most \(2k-1\). Part (a) gives the
    matching lower bound \(\min(2k-1,m)=2k-1\). These examples therefore
    have exactly the required \(k\)-rank, with \(m>2k_u-1\).

    \item[(c)] Assume first that \(n\geq2\) and that the factors have
    \(k\)-rank \(n\); that is, every \(n\) of the \(a_i\) are
    linearly independent. This condition holds for generic factors.
    Write
    \[
    r=2n-1,\qquad b_i=a_i\odot a_i,\qquad
    B=[b_1\ \cdots\ b_r]\in\mathbb R^{n^2\times r}.
    \]
    By part (a), the \(k\)-rank of \(B\) is at least
    \(\min(2n-1,r)=r\). Thus \(B\) has full column rank.

    \textbf{Step 1: Group the modes.}
    Group modes \((1,2)\), \((3,4)\), and \(5\), using the
    column-by-column ordering of \(\operatorname{vec}\). This reshapes
    \(T\) into the third-order tensor
    \[
    G=\sum_{i=1}^r b_i\otimes b_i\otimes a_i
    \in\mathbb R^{n^2\times n^2\times n}.
    \]
    This is a third-order decomposition with two full-column-rank
    factor matrices, so it can be recovered by simultaneous
    diagonalization as follows.

    \textbf{Step 2: Recover the directions of the \(b_i\).}
    Choose independent standard Gaussian vectors \(x,y\in\mathbb R^n\).
    Contract the last mode to form
    \[
    M_z=G(\cdot,\cdot,z)
       =B D_z B^\top,\qquad
    D_z=\operatorname{diag}(a_1^\top z,\ldots,a_r^\top z),
    \quad z\in\{x,y\}.
    \]
    With probability one, all \(a_i^\top y\) are nonzero and the
    ratios
    \[
    \rho_i=\frac{a_i^\top x}{a_i^\top y},\qquad 1\leq i\leq r,
    \]
    are pairwise distinct. Indeed, the \(k\)-rank assumption implies
    that every \(a_i\) is nonzero and no two are proportional.
    Each forbidden equality is therefore the zero set of a nonzero
    polynomial in \((x,y)\), and there are only finitely many such
    equalities.

    Compute an orthonormal basis
    \(Q\in\mathbb R^{n^2\times r}\) for
    \(\operatorname{range}(M_y)=\operatorname{range}(B)\).
    Let \(C=Q^\top B\), which is invertible. The computable
    \(r\times r\) matrix
    \begin{align*}
    K&=(Q^\top M_xQ)(Q^\top M_yQ)^{-1}\\
     &=(C D_x C^\top)(C D_y C^\top)^{-1}
      =C\operatorname{diag}(\rho_1,\ldots,\rho_r)C^{-1}
    \end{align*}
    has distinct eigenvalues. Its eigenvectors are the columns of
    \(C\), up to nonzero scaling and permutation. Multiplying these
    eigenvectors by \(Q\) therefore recovers vectors
    \(\widehat b_i=s_i b_i\), with \(s_i\ne0\), in some order.
    If \(M_y\) has rank less than \(r\) or \(K\) has repeated
    eigenvalues, resample \(x,y\).

    \textbf{Step 3: Recover the factors, including their scales.}
    Reshape each \(\widehat b_i\) into an \(n\times n\) matrix:
    \[
    \operatorname{unvec}(\widehat b_i)=s_i a_i a_i^\top.
    \]
    Any nonzero column is a vector \(d_i=t_i a_i\) for some
    nonzero real \(t_i\). Solve the linear system in the \(r\)
    scalar unknowns \(\lambda_i\):
    \[
    T=\sum_{i=1}^r\lambda_i d_i^{\otimes5}.
    \]
    A solution exists, namely \(\lambda_i=t_i^{-5}\). It is unique:
    the vectors \(d_i^{\otimes2}=t_i^2b_i\) are independent, so there
    are linear functionals \(L_j\) with
    \(L_j(d_i^{\otimes2})=\delta_{ji}\). Applying \(L_j\) to the
    first two modes of a relation
    \(\sum_i\gamma_i d_i^{\otimes5}=0\) gives
    \(\gamma_j d_j^{\otimes3}=0\), forcing \(\gamma_j=0\).
    Finally output
    \[
    \widehat a_i=\sqrt[5]{\lambda_i}\,d_i,
    \]
    taking the unique real fifth root. Since
    \(\sqrt[5]{t_i^{-5}}=t_i^{-1}\), this recovers every original
    \(a_i\) exactly, up to permutation. Thus the algorithm succeeds
    with probability one in exact arithmetic under the stated
    \(k\)-rank condition, even though \(r=2n-1>n\).

    If \(n=1\), then \(r=1\) and the tensor has a single entry;
    its real fifth root is the required scalar factor.
\end{enumerate}
\newpage

\end{document}
